\begin{statement}{A.5.7}{Proved}{A pure double conductor forces a smooth cusp image}
Let $S$ be a reduced Gorenstein complex surface and
$\pi:\overline S\to S$ its finite normalization.
Suppose that, over a neighbourhood of a point $p\in S$,
$\overline S$ is smooth and its conductor ideal is exactly
$\mathcal O_{\overline S}(-2C)$ for a smooth reduced
divisor $C$; its local branches over $p$ are disjoint.
Suppose that the generic transversal of $S$ along the
reduced image $D=\pi(C)$ is an ordinary cusp.
Then $D$ is smooth at $p$, $C\to D$ is an isomorphism
near the preimage of $p$, and the germ of $S$ at $p$ is
analytically a product of an ordinary cusp and a smooth disc.
No complete-intersection hypothesis on $D$ is required.
\end{statement}
\begin{proof}
Write $R=\mathcal O_{S,p}$ and let $\overline R$ be its
finite semilocal normalization. It is a product of regular
two-dimensional local analytic rings. Write $\mathfrak c$
for the conductor, and set
\[
 A=R/\mathfrak c,\qquad
 \overline A=\overline R/\mathfrak c,\qquad
 Q=\overline R/R=\overline A/A.
\]
We first give the exact duality comparison needed below.
The $R$-module $\overline R$ is maximal Cohen--Macaulay.
Since $R$ is Gorenstein,
$\operatorname{Ext}_R^j(\overline R,R)=0$ for $j>0$.
Its dual $\operatorname{Hom}_R(\overline R,R)$ is
precisely the conductor, viewed as a fractional ideal:
a homomorphism is multiplication in the total quotient ring,
and multiplication takes $\overline R$ into $R$ exactly
for conductor elements. Applying $\operatorname{Hom}_R(-,R)$
to $0\to R\to\overline R\to Q\to0$ therefore gives
\[
 \operatorname{Ext}_R^1(Q,R)=A.                    \tag{A.J.11}
\]
The depth lemma shows that $Q$ is a Cohen--Macaulay module
of dimension one. Gorenstein duality on such modules is the
involutive functor $M\mapsto\operatorname{Ext}_R^1(M,R)$.
Consequently $A$ is Cohen--Macaulay and
\[
 Q\simeq\operatorname{Ext}_R^1(A,R)=\omega_A,
 \qquad \operatorname{End}_A(Q)=A.                \tag{A.J.12}
\]
The final equality is the canonical multiplication map:
duality identifies its endomorphism ring with
$\operatorname{End}_A(A)=A$.

The ordinary-cusp hypothesis says that $A$ is generically
reduced. Its nilradical is therefore supported at finitely
many points; a one-dimensional Cohen--Macaulay ring has
no nonzero submodule with that support. Hence $A$ is reduced.
Choose local coordinates $(s_i,z_i)$ on the normalization
branches with $C=(z_i=0)$. Put
\[
 B=\prod_i\mathbb C\{s_i\},\qquad
 K=\prod_i\operatorname{Frac}\bigl(\mathbb C\{s_i\}\bigr).
\]
The specified conductor gives the actual ring comparison
\[
 \overline A=B[\epsilon]/(\epsilon^2),             \tag{A.J.13}
\]
component by component. The projection
$\overline A\to B$ is injective on $A$, since its kernel
is nilpotent and $A$ is reduced. It identifies $B$ with
the finite normalization of the curve ring $A$.

Under (A.J.13) the inclusion of $A$ takes the form
\[
 a\longmapsto a+\delta(a)\epsilon,
\]
where $\delta:A\to B$ is a complex-linear analytic
derivation. It extends to the total quotient ring, and on
each normalization branch has the form
$\delta(a)=u_i\,da/ds_i$ for a meromorphic germ $u_i$.
Let $u=(u_i)$ and write derivatives componentwise.
The map
\[
 \Phi:B[\epsilon]/(\epsilon^2)\longrightarrow K,
 \qquad b+c\epsilon\longmapsto c-u b'
                                                        \tag{A.J.14}
\]
is $A$-linear for this specified action: multiplication by
$a+u a'\epsilon$ changes its value to
$ac+u a'b-u(ab)'=a(c-u b')$. Its kernel at every generic
point is exactly the graph of the extended derivation on
$K$. It therefore induces a generic isomorphism from $Q$
to its image in $K$. The module $Q\simeq\omega_A$ is
torsion-free on the reduced curve, so this induced map is
injective.

Differentiation of convergent power-series germs onto
convergent power-series germs is surjective, by termwise
integration. Hence the image in (A.J.14) is exactly
\[
 I=B+uB\subset K.                                \tag{A.J.15}
\]
It is an invertible fractional ideal on the semilocal
normalization $B$. Every $A$-linear endomorphism of this
rank-one module is multiplication by an element of $K$;
preserving $I$ is equivalent, branch by branch, to belonging
to $B$. Thus $\operatorname{End}_A(I)=B$ as subrings of
$K$. Combining this with the actual $A$-linear isomorphism
$Q\xrightarrow{\Phi}I$ and (A.J.12) gives $A=B$.
Because $A$ is local, there is a single branch over $p$,
and its ring is regular. This proves that $D$ is smooth
and that $C\to D$ is an isomorphism there.

Finally $A=B=\mathbb C\{s\}$ makes $u=\delta(s)$
holomorphic. The analytic coordinate change
$s'=s+u(s)z$ has invertible differential at $z=0$ and
identifies the image of $A$ modulo $z^2$ with the functions
of $s'$ alone. Since $R$ is the inverse image of $A$ under
$\overline R\to\overline A$, it follows that
\[
 R=\mathbb C\{s'\}+z^2\mathbb C\{s',z\}
   =\mathbb C\{s',z^2,z^3\}
   \simeq\mathbb C\{s',x,y\}/(y^2-x^3).
                                                        \tag{A.J.16}
\]
This is the asserted product germ.
\end{proof}


% A.5.8 from part_A6_critical_curves.tex

